% Insert within A.137 after its finite geometric quotient construction.
% This fragment proves the replacement Gamma_136 -> Gamma_18, and exact class17.
\paragraph{The exact nilpotency class of the actual finite quotient.}
The geometric group $Q=\Pi/K$ above has nilpotency
class exactly $17$.  Thus its actual lower-central
quotient map can be sharpened to
\[
                            \Pi/\Gamma_{18}\Pi\twoheadrightarrow Q.
\]
Here is the covering-graph calculation, including a
specific nonzero commutator of weight seventeen.

Choose the free generators $x,y$ of $\Pi$ so that
the index-two character has values $0,1$ on them.
The Schreier generators of $\Lambda$ are
$a=x$, $b=yxy^{-1}$ and $z=y^2$.
After passing to $\Lambda^{\rm ab}/4$ they commute
and have order four.
Conjugation by $y$ exchanges $a,b$ and fixes $z$.
Consequently the exact group $P=\Pi/H$ is
\[
 \begin{gathered}
 P=(C_4\times C_4)\rtimes C_8,\qquad |P|=128,\\
 (a,b,k)(a',b',k')
   =\bigl((a,b)+s^k(a',b'),\,k+k'\bigr),\\
       s(a',b')=(b',a'),\quad
       a,b,a',b'\in\Z/4,\quad k,k'\in\Z/8,\\
                       x=(1,0,0),\qquad y=(0,0,1).
 \end{gathered}
\]
This presentation has the actual subgroup
$\langle a,b,y^2\rangle=C_4^3$ and the indicated
index-two quotient, so its kernel is exactly $H$.
With $[u,v]=uvu^{-1}v^{-1}$ one has
$[x,y]=ab^{-1}$,
$[[x,y],y]=(ab^{-1})^2\ne1$,
and every weight-four commutator in $P$ is trivial.
Thus $P$ has class three.

Take the covering graph of the two-loop rose with
vertices $g\in P$ and directed edges
$e_{g,x}:g\to gx$ and $e_{g,y}:g\to gy$.
Over $\F_2$ set
\[
 \begin{gathered}
 C_1=\bigoplus_{g\in P}\bigl(\F_2e_{g,x}\oplus\F_2e_{g,y}\bigr),
       \qquad C_0=\bigoplus_{g\in P}\F_2\delta_g,\\
 \partial e_{g,x}=\delta_g+\delta_{gx},\qquad
 \partial e_{g,y}=\delta_g+\delta_{gy},\qquad
             h\cdot e_{g,u}=e_{hg,u}.
 \end{gathered}
\]
Left translation commutes with the boundary map.
The following is an explicit finite group:
\[
 \begin{gathered}
 \mathcal Q=\{(v,g)\in C_1\times P:
                          \partial v=\delta_g+\delta_1\},\\
 (v,g)(w,h)=(v+g\cdot w,gh),\qquad
                  (v,g)^{-1}=(g^{-1}\cdot v,g^{-1}),\\
                   X=(e_{1,x},x),\qquad Y=(e_{1,y},y).
 \end{gathered}
\]
The free-generator map $\Pi\to\mathcal Q$ is
the actual path map modulo two.
For an element of $H$ it records a closed path
in this graph.
A spanning tree gives $256-128+1=129$ independent
non-tree-edge cycles; the corresponding Schreier
loops form the free generators of $H$.
Hence its restriction induces the exact isomorphism
\[
        H/\overline{H^2[H,H]}\ \xrightarrow{\ \sim\ }
                       Z_1:=\ker(\partial:C_1\to C_0).
\]
For each endpoint $g$, choose a word reaching it.
Adding any cycle by an element of $H$ gives every
pair in the displayed fibre over $g$.
Thus $\Pi\to\mathcal Q$ is onto, its kernel is
exactly $K$, and it identifies $Q$ with this
specific group.
In particular $|Q|=128\cdot2^{129}=2^{136}$.
All these maps have finite targets and are
continuous for the profinite groups in A.137.

This gives the actual comparison with the established
relation-module construction.  In
\cite[Definition 2.1 and Theorem 3.7]{AgentDietrichHulpke2020},
take the finite group to be $P$, the free rank to be two,
and the prime to be two.  The map sending the path image
of a word $w$ to its image $\Psi_2(w)$ in that theorem
is a well-defined group isomorphism: the two kernels
are exactly $\overline{H^2[H,H]}$, as proved here and
in that theorem.  On the kernel, the edge maps
$e_{g,x}\mapsto(g,0)$ and $e_{g,y}\mapsto(0,g)$
identify our boundary sequence with the mod-two
relation sequence in
\cite[Proposition 2.8]{AgentYamin1991}; its displayed
boundary is right multiplication by $x-1,y-1$ on
the two left regular summands.
The specific lower-central calculation follows below.

Since $P$ has class three, $\Gamma_4Q$ is a
subspace of the elementary abelian cycle group
$Z_1$.
Its exact generating vectors are the $P$-translates
of the four actual commutators
\[
 [[[X,Y],X],X],\quad [[[X,Y],X],Y],\quad
 [[[X,Y],Y],X],\quad [[[X,Y],Y],Y].
\]
To justify completeness, the normal closure of
$[X,Y]$ is $\Gamma_2Q$.
If a normal subgroup is normally generated by
a list, its commutator with $\langle X,Y\rangle$
is normally generated by the commutators of
that list with $X,Y$: in the quotient killing
those latter commutators every original
generator is central.
Apply this twice.
On $Z_1$ conjugation depends only on $P$,
giving precisely the stated vector span.

For a $P$-stable subspace $W\subset Z_1$ the next
commutator subspace is exactly
\[
                  [W,Q]=(x-1)W+(y-1)W.
\]
Indeed $[(v,1),X]=(v+x\cdot v,1)$, and similarly
for $Y$.
The identity
$(gh-1)w=(g-1)(h\cdot w)+(h-1)w$
expresses the difference for every word in $x,y$
through these two images.
It also proves their sum is $P$-stable.

The retained integer-bitset elimination uses
the indexing
$g=a+4b+16k$ and $e_{g,x}=2g$, $e_{g,y}=2g+1$.
It starts with those four literal commutators
and their $128$ translates, and applies the
last displayed recurrence until zero.
Its exact dimensions are
\[
 \begin{gathered}
 \begin{array}{c|rrrrrrrr}
 j&4&5&6&7&8&9&10&11\\ \hline
 \dim_{\F_2}\Gamma_jQ&124&121&115&107&97&85&71&55
 \end{array}
 \\[6pt]
 \begin{array}{c|rrrrrrr}
 j&12&13&14&15&16&17&18\\ \hline
 \dim_{\F_2}\Gamma_jQ&41&29&19&11&5&2&0 .
 \end{array}
 \end{gathered}
\]
Each retained basis element is an unreduced
independent original generator rather than
an unspecified linear combination.
Its recorded provenance therefore remains
an actual commutator at every step.

Explicitly, start with $C_4=[[[X,Y],X],X]$
and successively replace $C$ by $[C,Z]$
for the thirteen letters
\[
                         Z:\quad \mathrm{XXYXYXYYYYYYY}.
\]
The resulting weight-seventeen element is
\[
                    \left(\sum_{g\in P}e_{g,x},\,1\right)\ne1.
\]
Its bitset is the $256$-bit integer with every
even-indexed bit equal to one and every
odd-indexed bit zero.
It is nonzero, is a cycle, and is fixed by
both left translations.
The final zero row proves $\Gamma_{18}Q=1$;
this explicit element proves $\Gamma_{17}Q\ne1$.
Together they prove the asserted exact class and
the sharpened quotient map.

The code
\path{code/harmonic_quotient_class_certificate.py},
its full output, and
\path{code/harmonic_quotient_class_data.json}
retain every basis vector and its commutator
provenance.
The computation also checks all $128^3$
associative products of $P$, every translated
edge boundary, every canonical word lift,
the four starting paths, and the equality of
the recorded deepest word with its displayed
cycle.  No enumeration of the $2^{136}$
elements of $Q$ is used.

\paragraph{An independent augmentation-ideal upper bound.}
There is a second exact verification of the upper bound.
Put $\mathcal R=\F_2[P]$ and let $I$ be its augmentation
ideal.  The independent program
\path{code/harmonic_quotient_augmentation_audit.py}
computes its powers in the regular module by
$v\mapsto v+xv$ and $v\mapsto v+yv$.
The identity for $gh-1$ above proves that their span
on $I^j$ is exactly $I^{j+1}$.
Its dimensions are
\[
 \begin{aligned}
 (\dim I^j)_{j=0}^{17}
   ={}&(128,127,125,121,115,107,97,85,71,\\
       &57,43,31,21,13,7,3,1,0).
 \end{aligned}
\]
The successive differences agree coefficient by
coefficient with $(1+t)^2(1+t^2)^3(1+t^4)^2$;
in particular the executed finite calculation proves
$I^{17}=0$ and $\dim I^{16}=1$ directly.

For a free word $w$, write its path image as
$(D(w),\nu(w))$.  The multiplication already proved gives
$D(uv)=D(u)+\nu(u)D(v)$.
With $g=\nu(u)$ and $h=\nu(v)$, the exact identities are
\[
 \begin{aligned}
 [g,h]-1
   &\;=\bigl((g-1)(h-1)-(h-1)(g-1)\bigr)g^{-1}h^{-1},\\
 D([u,v])
   &\;=(1-ghg^{-1})D(u)+(g-[g,h])D(v),\\
 1-ghg^{-1}&\;=-g(h-1)g^{-1},\qquad
 g-[g,h]=gh(1-g^{-1})h^{-1}.
 \end{aligned}
\]
They prove simultaneously by induction that
\[
 \nu(w)-1\in I^j,\qquad D(w)\in (I^{j-1})^2
                              \quad(w\in\Gamma_j\Pi).
\]
Indeed commutators of weights $a,b$ have endpoint
difference in $I^{a+b}$ and chain in $(I^{a+b-1})^2$.
Products and inverses preserve these ideals.
Normal conjugation also preserves them by the literal formula
\[
 D(vuv^{-1})
   =(1-\nu(v)\nu(u)\nu(v)^{-1})D(v)+\nu(v)D(u).
\]
The finite target makes the assertion valid on
the closed profinite lower-central subgroups.
For $w\in\Gamma_{18}\Pi$, its endpoint is one
because $P$ has class three, and its chain is zero
because $I^{17}=0$.
The exact path kernel therefore gives
$\Gamma_{18}\Pi\subset K$ independently of the
preceding lower-central row-space calculation.
The same independent program expands the specified
weight-seventeen word into $196606$ letters and
traces its edges one by one, obtaining all $128$
$x$-edges and no $y$-edges.
This separately verifies its nonzero image.
The exact class of this detecting quotient does
not assert that the constructed pair first fails
at class seventeen.
